\documentclass[11pt,a4paper]{article}
\usepackage[utf8]{inputenc}
\usepackage[T1]{fontenc}
\usepackage[brazilian]{babel}
\usepackage{geometry}
\usepackage{enumitem}
\usepackage{booktabs}
\usepackage{amsmath}
\usepackage{graphicx}
\graphicspath{{../figures_global/}}
\usepackage{fancyhdr}
\usepackage{hyperref}

\geometry{margin=2.5cm, headheight=32pt}
\setlength{\parindent}{0pt}
\setlength{\parskip}{0.5em}

\pagestyle{fancy}
\fancyhf{}
\fancyhead[L]{\small\textbf{Faculdade de Ciências e Tecnologias em Engenharia}}
\fancyhead[R]{\includegraphics[height=0.85cm]{../figures_global/unb.png}}
\renewcommand{\headrulewidth}{0.4pt}
\fancyfoot[C]{\small Laboratório de Comunicações --- Práticas}
\fancyfoot[R]{\small \thepage}
\renewcommand{\footrulewidth}{0.4pt}
\fancypagestyle{plain}{%
  \fancyhf{}%
  \fancyhead[L]{\small\textbf{Faculdade de Ciências e Tecnologias em Engenharia}}%
  \fancyhead[R]{\includegraphics[height=0.85cm]{../figures_global/unb.png}}%
  \renewcommand{\headrulewidth}{0.4pt}%
  \fancyfoot[C]{\small Laboratório de Comunicações --- Práticas}%
  \fancyfoot[R]{\small \thepage}%
  \renewcommand{\footrulewidth}{0.4pt}%
}

\newcommand{\pratica}[3]{%
  \subsection{#1}
  \textit{Objetivo.} #2

  \textit{Método e medições.} #3
  \vspace{0.3em}
}

\title{\textbf{Catálogo de Práticas}\\[0.3em]\large Laboratório de Comunicações}
\author{Faculdade de Ciências e Tecnologias em Engenharia --- FCTE/UnB}
\date{}

\begin{document}

\thispagestyle{plain}
\maketitle
\vspace{-1.0em}

\section{Apresentação}

Este documento propõe as práticas do Laboratório de Comunicações, organizadas em duas frentes complementares. A Frente A trata de modulação, recepção e sincronização em rádio definido por software com o ambiente GNU Radio. A Frente B trata de radiopropagação e caracterização de circuitos analógicos de radiofrequência e de seu impacto sobre o desempenho dos sistemas de comunicação. Cada prática indica seu objetivo e o conjunto de medições e análises esperadas.

\section{Carga horária}

O componente tem 2 créditos, distribuídos entre atividades teórico-práticas e extensão. As oito práticas deste catálogo ocupam a parte teórico-prática, com quatro práticas em cada frente. A componente de extensão é dedicada ao apoio técnico às equipes de competição, quando as práticas pertinentes são retomadas e adaptadas ao projeto real da equipe.

\begin{center}
\small
\begin{tabular}{@{}p{0.60\textwidth} p{0.24\textwidth}@{}}
\toprule
\textbf{Atividade} & \textbf{Práticas} \\
\midrule
Frente A --- GNU Radio & 4 \\
Frente B --- Radiofrequência & 4 \\
Extensão --- apoio às equipes & --- \\
\midrule
\textbf{Total} & \textbf{8} \\
\bottomrule
\end{tabular}
\end{center}

\section{Frente A --- Modulação, Recepção e Sincronização em GNU Radio}

\pratica{Introdução ao SDR e análise espectral}
{Familiarizar o estudante com a cadeia de recepção em rádio definido por software e com a medição de espectro no GNU Radio.}
{Montagem de um fluxograma de recepção, ajuste de frequência, taxa de amostragem e ganho, estimação de espectro por FFT com janelamento, e captura de um sinal do ambiente.}

\pratica{Modulação e demodulação analógica AM e FM}
{Implementar e avaliar as modulações analógicas em amplitude e em frequência.}
{Geração de AM e FM com parâmetros ajustáveis, análise espectral, verificação da regra de Carson e demodulação por detector de envoltória, coerente e discriminador.}

\pratica{Modulação digital e recepção com filtro casado}
{Gerar sinais digitais passa-banda e avaliar a recepção ótima com filtro casado.}
{Geração de BPSK e QPSK, conformação por cosseno levantado, diagrama de olho e constelação, filtro casado e medição da taxa de erro em função da relação sinal-ruído.}

\pratica{Sincronização e enlace digital ponta a ponta}
{Integrar sincronização de portadora, de temporização e de quadro em um enlace completo.}
{Malha de captura de fase e de Costas, recuperação de temporização, detecção de preâmbulo, transmissão de uma mensagem entre dois rádios e comparação da taxa de erro com a curva teórica.}

\section{Frente B --- Radiopropagação e Caracterização de Circuitos Analógicos}

\pratica{Linhas de transmissão e casamento de impedância}
{Compreender reflexão, coeficiente de onda estacionária e casamento de impedância.}
{Medição de parâmetros de espalhamento no analisador de redes, uso da carta de Smith, e projeto e verificação de uma rede de casamento.}

\pratica{Radiopropagação e diagrama de antenas}
{Verificar modelos de propagação e levantar as características de radiação de antenas.}
{Medição da potência recebida em função da distância, ajuste do expoente de perda de percurso, e levantamento do diagrama de radiação e do ganho de uma antena.}

\pratica{Amplificador de baixo ruído e linearidade}
{Caracterizar ganho, ruído e linearidade de dispositivos ativos de radiofrequência.}
{Medição de ganho e figura de ruído, cálculo do efeito sobre a sensibilidade pela fórmula de Friis, ponto de compressão de um decibel e ensaio de dois tons para intermodulação.}

\pratica{Orçamento de enlace e impacto das imperfeições}
{Consolidar os parâmetros de um enlace e relacionar imperfeições de circuito ao desempenho.}
{Montagem de um enlace curto, cálculo do orçamento com ganhos, perdas e ruído, e observação do efeito do ruído de fase do oscilador sobre a constelação e o vetor de erro.}

\section{Articulação com a componente de extensão}

As práticas das duas frentes fornecem a base técnica para o apoio às equipes de competição, previsto na componente de extensão. A partir de uma demanda concreta da equipe, como um enlace de telemetria de foguete ou de minissatélite, a turma seleciona e adapta as práticas pertinentes, executa as medições no contexto real do projeto da equipe e entrega um laudo técnico com recomendações. Assim, o mesmo instrumental usado nas práticas didáticas passa a atender a uma necessidade externa, cumprindo o caráter extensionista do componente.

\end{document}
